% SSSUP-FONTE-STORICA: fonti-locali/originali/2014-2020.zip
% Anno 2020, esercizio 3. Testo storico; figura ricostruita da PDF.
% Verificare editorialmente la trasposizione prima di considerarla definitiva.
\renewcommand{\Anno}{2020}
\renewcommand{\NumeroProblema}{3}

\renewcommand{\Testo}{
Un'aula deve essere riconfigurata per il rientro durante una pandemia e si decide che ogni studente deve distare da ogni altro almeno 2 metri. Si considera ogni studente come un punto, da posizionare opportunamente nell'aula, incluso sulle linee che rappresentano i muri.
\begin{enumerate}
\item[(a)] È possibile mettere in un'aula quadrata di lato $2n$ metri più di $(n+1)^2$ studenti? Se sì\, per quali $n$?
\item[(b)] Per la cattedra e la lavagna è necessario lasciare una striscia verso i banchi, di almeno 2 metri sul corrispondente lato dell'aula (il lato  e la larghezza della striscia possono essere scelti). Tra le aule rettangolari di area 36 $m^2$ e con lati di almeno 4 metri, quale contiene più studenti? Quanti studenti ci si possono sedere e con che configurazione? Quanto sarà larga la striscia?
\end{enumerate}
}

\renewcommand{\Risposta}{}
\renewcommand{\Aiutino}{}
\renewcommand{\SoluzioneUfficiale}{}
\renewcommand{\SoluzioneRedazionale}{}
\renewcommand{\SoluzioniAlternative}{}
\renewcommand{\Osservazioni}{}
\renewcommand{\Commenti}{}
